\documentclass[10pt,a4paper]{amsart}
\usepackage[margin=19mm]{geometry}
\usepackage{amsmath,amssymb,mathtools}
\usepackage[scr=rsfs,cal=euler]{mathalfa}
\usepackage{xfrac}
\usepackage[unicode,hidelinks,bookmarks=false]{hyperref}
\newcommand{\ie}{i.e.\ }
\newcommand{\cf}{cf.\ }
\newcommand{\resp}{respectively\ }
\newcommand{\Subsection}{\subsection}
\providecommand{\nobreakdash}{\nobreak\hbox{-}}
\setlength{\emergencystretch}{2em}
\multlinegap0pt
\usepackage{xcolor}
\usepackage{cancel}
\usepackage{amssymb}
\usepackage{dsfont}
\usepackage{mathtools}
\usepackage{algorithm}\usepackage{algpseudocode}
\usepackage[shortlabels]{enumitem}
\usepackage{mdframed}
\usepackage[normalem]{ulem}
\usepackage{bm}
\usepackage{tcolorbox}
\newtheorem{claim}{Claim}
\def\bigabs#1{\left| #1 \right|}
\def\gI{{\mathcal{I}}}
\def\gQ{{\mathcal{Q}}}
\def\gS{{\mathcal{S}}}
\title{Near-Optimal Clustering in Mixture of Markov Chains}
\author{Junghyun Lee, Yassir Jedra, Alexandre Prouti\`ere, Se-Young Yun}
\date{Source: arXiv:2506.01324v3 (CC BY 4.0)}
\begin{document}
\maketitle

\noindent\emph{Source and licence.} Near-Optimal Clustering in Mixture of Markov Chains. Version arXiv:2506.01324v3 (CC BY 4.0), distributed by arXiv under the Creative Commons Attribution 4.0 International licence, http://creativecommons.org/licenses/by/4.0/, as recorded in the arXiv licence metadata for this exact version and shown on the abstract page https://arxiv.org/abs/2506.01324v3. Full licence notice: \texttt{licenses/arxiv-2506.01324-CC-BY-4.0.txt}.\par\smallskip

\noindent\textit{Editorial scope.} Counted unit: the article's claim claim (source0.tex lines 2318--2320). The article's complete original proof of it is delivered with it (source0.tex lines 2321--2340, one \texttt{proof} environment of 20 lines). No mathematics is written, repaired or re-derived: the statement and the proof are the article's own lines, in the article's own notation, and no retained line is substituted or re-flowed.  No further source-local context is needed: every cross-reference of every retained line resolves inside the two delivered spans.  Nothing was omitted from any retained span, so the package carries no omission record.  No retained line contains a citation, so the wrapper carries no bibliography and no bibliography entry.  No retained line contains a figure environment, an image-inclusion command or drawing code, so no image is delivered and the package declares no asset dependency.  Editorial formatting only, in addition: an AMS \texttt{amsart} preamble that restates the article's own declarations and macros used by the retained lines, namely the environments \texttt{claim} on separate counters, together with the standard packages those lines need.  The retained environments print in this document as Claim 1 in order of appearance, whereas the article prints the same items with the numbering of its own full document.  The complete native source of the article as distributed in the arXiv e-print for this exact version is bundled unchanged as \texttt{sources/179149/source0.tex}.
\begin{claim}
    $\widehat{K} = K$.
\end{claim}

\begin{proof}
    Without loss of generality, assume that $|\gI_1| \geq |\gI_2| \geq \cdots \geq |\gI_K|$.
    
    We first show that for each $k \in [K]$,
    \begin{equation*}
        \exists t_k^\star \in \bigcup_{k' \in [K]} \gI_{k'} \setminus \bigcup_{\ell=1}^{k-1} \gS_\ell \quad \text{s.t.} \quad \bigabs{\gQ_{t_k^\star} \setminus \bigcup_{\ell=0}^{k-1} \gS_\ell} \geq |\gI_k|.
    \end{equation*}
    Due to the properties (iii) and (iv), and the greedy nature of lines 7-10, the above is indeed true.
    
    Let $\{t_1^\star, \cdots, t_K^\star\}$ be the selected ``centers'' with $\gI_k \subseteq \gQ_{t_k^\star}$.

    Now, we show that after $k$ has reached $K + 1$ in line 10, the \textbf{while} loop terminates.
    By property (ii), the number of remaining trajectories is
    \begin{equation*}
        \bigabs{[T] \setminus \bigcup_{k \in [K]} \gQ_{t_k^\star}}
        \leq \bigabs{[T] \setminus \bigcup_{k \in [K]} \gI_k}
        \overset{(ii)}{\leq} \frac{32 R T}{\log\frac{TH}{\delta}},
    \end{equation*}
    which is precisely the termination criterion at line 7.
\end{proof}

\end{document}
